\chapter[Arquitetura de Software]{Arquitetura de Software}
\label{cap:software}

Este capítulo descreve o software do robô: a base do sistema operacional e
das bibliotecas de acesso ao hardware, as ferramentas de teste
desenvolvidas, a organização em nós do ROS~2 e a forma como os dados
fluem entre os sensores, os algoritmos de mapeamento e navegação e os
motores.

\section{Visão geral}

O software está dividido entre dois computadores que se comunicam por
Wi-Fi. O critério da divisão é a capacidade de processamento. O Raspberry
Pi 3 embarcado executa apenas o que precisa estar perto do hardware: os
drivers dos sensores e o controle dos motores. O notebook, que funciona
como estação de base, executa a fusão de sensores, o mapeamento, a
navegação e a visualização, que exigem mais processador e memória. As duas
partes usam a mesma distribuição do ROS~2, a Jazzy, o que garante que as
mensagens trocadas sejam compatíveis.

O desenvolvimento seguiu duas fases. Na primeira, cada componente foi
validado isoladamente com programas próprios em Python e C, sem ROS
(Seção~\ref{sec:bancada}). Na segunda, os componentes validados passam a
ser integrados como nós do ROS~2 (Seções~\ref{sec:ros} a
\ref{sec:integracao-sw}).

\section{Sistema operacional e acesso ao hardware}

O Raspberry Pi executa o Raspberry Pi OS ``Trixie'' de 64 bits, com
Python 3.13. A interface gráfica foi desativada
(\texttt{multi-user.target}) para liberar memória, e o acesso é feito por
SSH a partir do notebook.

O acesso aos pinos de GPIO é feito pela biblioteca pigpio \cite{pigpio}.
Ela foi escolhida por três motivos. Ela dá acesso aos periféricos de PWM
por hardware, necessários para acionar os MD10C com frequência estável.
Ela amostra os pinos de entrada por DMA e entrega cada borda com o
instante em que ocorreu, o que permite contar os pulsos dos encoders sem
depender de interrupções do sistema operacional. E ela funciona como um
serviço (\texttt{pigpiod}) que atende vários programas ao mesmo tempo. Uma
consequência desse modelo é que o PWM continua ativo se o programa que o
configurou terminar de forma anormal. Por isso, todos os programas que
acionam motores zeram o PWM ao sair, e existe um comando de emergência
para isso. Como o Trixie não traz o pacote da pigpio, ela foi compilada a
partir do código-fonte.

\section{Ferramentas de validação dos componentes}
\label{sec:bancada}

Antes da integração com o ROS~2, cada componente foi testado com um
programa próprio. Isso permite separar defeitos de montagem de defeitos de
software.

\begin{itemize}
	\item \textbf{\texttt{teste\_motores.py}}: bancada de teste dos
	motores e encoders, executada com as rodas suspensas. Ela tem um modo
	para cada etapa da caracterização. O modo \texttt{encoder} mede as
	contagens por volta da roda (Eq.~\ref{eq:cpr}) e confere o sentido de
	contagem. O modo \texttt{direcao} confere se cada lado gira no sentido
	comandado e se o encoder concorda. O modo \texttt{varredura} aumenta o
	ciclo de trabalho em passos de 5\% para levantar a zona morta e a
	curva de velocidade em regime. O modo \texttt{degrau} aplica um degrau
	de tensão e amostra a velocidade a 100~Hz para estimar a constante de
	tempo dos motores. Todos os comandos passam por rampas, e o programa
	interrompe o teste se um motor recebe tensão alta e o encoder não
	registra movimento, o que indica roda travada ou encoder desligado;
	\item \textbf{\texttt{teste\_encoders.py}}: mostra ao vivo os níveis
	lógicos, o número de bordas de cada canal e a posição de cada
	encoder, para diagnosticar fios soltos ou trocados;
	\item \textbf{\texttt{controle\_wasd.py}}: teleoperação pelo teclado,
	via SSH. Por segurança, o robô para sozinho se as teclas deixam de
	chegar, se a conexão SSH cai ou se um motor trava;
	\item \textbf{\texttt{lidar.py}} e \textbf{\texttt{ver\_lidar.py}}:
	leitura do RPLIDAR C1 no Raspberry Pi e exibição das varreduras em
	gráfico polar no notebook, com os dados enviados pela conexão SSH;
	\item \textbf{\texttt{profundidade.py}} e
	\textbf{\texttt{astra\_depth.c}}: leitura da profundidade da Astra Pro
	e envio ao notebook como vídeo MJPEG colorido.
\end{itemize}

A leitura da câmera de profundidade exigiu uma solução específica. O
pacote \texttt{openni} do repositório do Python não é compatível com o
Python 3.12 ou mais recente. Depois de corrigido, ele ainda corrompia a
memória, porque as estruturas de dados declaradas no pacote não
correspondiam às da biblioteca da Orbbec. A solução foi escrever um
programa em C compilado diretamente contra os cabeçalhos do SDK OpenNI2 da
fabricante. Ele lê os quadros de profundidade e os repassa ao programa em
Python. Na integração com o ROS~2, esse programa será substituído pelo
driver da própria câmera e servirá apenas como ferramenta de teste.

\section{ROS 2}
\label{sec:ros}

O ROS~2 (\textit{Robot Operating System 2}) é um conjunto de bibliotecas
e ferramentas para software de robôs \cite{macenski2022ros2}. Um sistema
em ROS~2 é formado por \textbf{nós}, processos independentes que trocam
\textbf{mensagens} de tipos padronizados por meio de \textbf{tópicos} com
nome. Um nó que publica em um tópico não precisa saber quais nós o
assinam. Por isso, um driver de sensor pode ser trocado sem mudar os
algoritmos que consomem seus dados. A comunicação usa o padrão DDS, que
descobre automaticamente os nós presentes na rede local. Assim, nós que
rodam em computadores diferentes se comunicam como se estivessem no mesmo
computador. É esse recurso que permite dividir o sistema entre o
Raspberry Pi e o notebook.

O ROS~2 também mantém a árvore de transformações (TF), que registra a
posição e a orientação relativas entre os sistemas de coordenadas do
robô: o mapa, a origem da odometria, o centro do robô e cada sensor. Os
algoritmos usam essa árvore para converter, por exemplo, um ponto medido
pelo LiDAR para coordenadas do mapa.

O ROS~2 Jazzy tem suporte oficial ao Ubuntu 24.04, mas não ao Raspberry Pi
OS Trixie nem ao Arch Linux, que são os sistemas usados no projeto. No
Raspberry Pi, ele foi instalado a partir do repositório comunitário
\textit{rospian}, que distribui o Jazzy compilado para o Trixie arm64.
Pacotes ausentes nesse repositório, como o driver do LiDAR, são compilados
na própria placa com o \texttt{colcon}. No notebook, o ROS~2 roda em um
contêiner Ubuntu 24.04 gerenciado pelo \textit{distrobox}, que compartilha
a rede e a interface gráfica com o sistema principal.

\section{Nós embarcados no Raspberry Pi}

\subsection{Driver do LiDAR}

O nó \texttt{sllidar\_node}, do pacote oficial \texttt{sllidar\_ros2} da
Slamtec, lê o RPLIDAR C1 pela porta serial e publica cada volta no tópico
\texttt{/scan}, como mensagem \texttt{sensor\_msgs/LaserScan}. Essa
mensagem contém as distâncias medidas e os ângulos correspondentes.

\subsection{Nó da base}

Os drivers MD10C não têm controle de velocidade próprio. Eles só aplicam
a tensão comandada. Por isso, o controle de velocidade das rodas roda no
Raspberry Pi, em um nó próprio, chamado aqui de nó da base. Ele assina o
tópico \texttt{/cmd\_vel}, que contém a velocidade linear $v$ e a
velocidade angular $\omega$ desejadas para o robô, e calcula a velocidade
de cada lado pelo modelo cinemático diferencial \cite{siegwart2011}:
\begin{equation}
	\label{eq:cinematica-inversa}
	v_e = v - \omega \frac{L}{2}, \qquad v_d = v + \omega \frac{L}{2},
\end{equation}
em que $v_e$ e $v_d$ são as velocidades das rodas esquerda e direita e
$L$ é a distância entre os centros das rodas dos dois lados.
% TODO: medir L e citar o valor aqui.

Cada lado tem um controlador PID que compara a velocidade desejada com a
velocidade medida pelo encoder e calcula o ciclo de trabalho aplicado ao
MD10C:
\begin{equation}
	\label{eq:pid}
	u(t) = K_p\, e(t) + K_i \int_0^t e(\tau)\, d\tau + K_d \frac{de(t)}{dt},
\end{equation}
em que $e(t)$ é o erro de velocidade e $K_p$, $K_i$ e $K_d$ são os ganhos
proporcional, integral e derivativo. Os parâmetros medidos na bancada (zona
morta, curva de velocidade e constante de tempo) servem de ponto de
partida para esses ganhos e para uma parcela de compensação direta
(\textit{feedforward}).

No sentido contrário, o nó da base converte as contagens dos encoders em
velocidades das rodas (Eq.~\ref{eq:deslocamento}) e calcula a velocidade
do robô pela cinemática direta:
\begin{equation}
	\label{eq:cinematica-direta}
	v = \frac{v_d + v_e}{2}, \qquad \omega = \frac{v_d - v_e}{L}.
\end{equation}
Integrando essas velocidades ao longo do tempo, o nó estima a posição e a
orientação do robô e as publica no tópico \texttt{/odom}
(\texttt{nav\_msgs/Odometry}). Por segurança, o nó para os motores se
nenhum comando chegar em um intervalo curto, o que cobre falhas de rede e
travamentos na estação de base.

\subsection{Driver da IMU}

O driver \texttt{bno055} \cite{bno055ros} lê a IMU pela UART e publica a
orientação, a velocidade angular e a aceleração no tópico \texttt{/imu}
(\texttt{sensor\_msgs/Imu}).

\subsection{Transformações estáticas}

A posição de cada sensor em relação ao centro do robô
(\texttt{base\_link}) é fixa e é publicada como transformação estática a
partir de uma descrição do robô em URDF. Exemplos são a posição do LiDAR,
próxima do centro do chassi e a 10~cm do chão, e a posição da IMU.

\section{Nós executados na estação de base}

\subsection{Fusão de sensores}

O pacote \texttt{robot\_localization} implementa um filtro de Kalman
estendido (EKF) que combina várias fontes de odometria em uma estimativa
única do movimento \cite{moore2016ekf}. Em um robô \textit{skid-steer}, a
rotação calculada pela Eq.~(\ref{eq:cinematica-direta}) é imprecisa,
porque as rodas escorregam lateralmente nas curvas. Por isso, o filtro
usa dos encoders apenas a velocidade linear $v$ e usa do giroscópio da IMU
a velocidade angular $\omega$. Essa configuração é a mesma adotada no
robô HelMoRo \cite{helmoro}, que tem sensores semelhantes aos deste
trabalho. O filtro publica a transformação \texttt{odom}$\rightarrow$%
\texttt{base\_link}. Por isso, o nó da base não deve publicá-la, para que
a árvore de transformações não tenha duas fontes para o mesmo elo.

\subsection{Mapeamento}

O mapeamento usa o \texttt{slam\_toolbox} \cite{macenski2021slam}. Esse
pacote resolve o problema de localização e mapeamento simultâneos (SLAM):
ele constrói o mapa a partir das varreduras do LiDAR e, ao mesmo tempo,
estima a posição do robô nesse mapa. A estimativa da odometria serve de
ponto de partida para alinhar cada nova varredura com as anteriores.
Quando o robô volta a um lugar já visitado, o pacote reconhece o lugar e
corrige o erro acumulado no caminho (fechamento de laço). O resultado é
um mapa de ocupação 2D, uma grade com células livres, ocupadas ou
desconhecidas, com resolução de 5~cm. O mapa pode ser salvo e reutilizado
na navegação. O \texttt{slam\_toolbox} publica a transformação
\texttt{map}$\rightarrow$\texttt{odom}, que corrige a deriva da
odometria.

\subsection{Navegação}

A navegação autônoma usa o Nav2 \cite{macenski2020nav2}. A partir do mapa,
da posição estimada e de mapas de custo que representam os obstáculos e
uma margem de segurança em torno deles, o Nav2 planeja um caminho até o
ponto de destino. Em seguida, um controlador local gera os comandos de
velocidade em \texttt{/cmd\_vel} que fazem o robô seguir esse caminho e
desviar de obstáculos novos detectados pelo LiDAR. O destino é escolhido
pelo usuário no RViz, a ferramenta de visualização do ROS~2.

Os comandos de velocidade podem vir do Nav2 ou da teleoperação. O pacote
\texttt{twist\_mux} escolhe entre essas fontes dando prioridade à
teleoperação, para que o operador possa assumir o controle a qualquer
momento.

\section{Integração do software}
\label{sec:integracao-sw}

A Tabela~\ref{tab:topicos} resume os tópicos que ligam os nós. Os dados
fluem em um ciclo fechado. O LiDAR, os encoders e a IMU publicam suas
medidas no Raspberry Pi. O EKF, no notebook, combina a odometria das
rodas com o giroscópio. O \texttt{slam\_toolbox} usa essa estimativa e as
varreduras do LiDAR para construir o mapa e localizar o robô. O Nav2 usa
o mapa e a localização para gerar comandos de velocidade. Esses comandos
voltam ao Raspberry Pi, onde o nó da base os converte em tensão nos
motores.

\begin{table}[htb]
	\centering
	\caption{Principais tópicos do sistema}
	\label{tab:topicos}
	\begin{tabular}{|l|l|l|l|}
		\hline
		\textbf{Tópico} & \textbf{Tipo} & \textbf{Publicado por} & \textbf{Usado por} \\ \hline
		\texttt{/scan} & \texttt{LaserScan} & \texttt{sllidar\_node} (Pi) & SLAM, Nav2 \\ \hline
		\texttt{/odom} & \texttt{Odometry} & nó da base (Pi) & EKF \\ \hline
		\texttt{/imu} & \texttt{Imu} & \texttt{bno055} (Pi) & EKF \\ \hline
		\texttt{/odometry/filtered} & \texttt{Odometry} & EKF (notebook) & Nav2 \\ \hline
		\texttt{/map} & \texttt{OccupancyGrid} & \texttt{slam\_toolbox} (notebook) & Nav2, RViz \\ \hline
		\texttt{/cmd\_vel} & \texttt{Twist} & Nav2 ou teleoperação & nó da base (Pi) \\ \hline
	\end{tabular}
	\legend{Fonte: Autor.}
\end{table}

A árvore de transformações segue a cadeia
\texttt{map}$\rightarrow$\texttt{odom}$\rightarrow$\texttt{base\_link}%
$\rightarrow$\{\texttt{laser}, \texttt{imu}\}. Cada elo tem um único
responsável: o \texttt{slam\_toolbox}, o EKF e a descrição do robô,
respectivamente.

Como os nós estão em computadores diferentes, a rede faz parte do
sistema. Só passam pelo Wi-Fi as varreduras do LiDAR, que são pequenas,
as mensagens de odometria e IMU e os comandos de velocidade. A câmera de
profundidade terá resolução e taxa de quadros reduzidas para não saturar
o enlace. A descoberta automática do DDS usa \textit{multicast}, que
algumas redes institucionais bloqueiam. Nesse caso, a comunicação é feita
por um roteador dedicado ou por um servidor de descoberta (\textit{Discovery
Server}) do próprio DDS.

\section{Ajuste dos controladores por enxame de partículas}

Os ganhos dos controladores PID da Eq.~(\ref{eq:pid}) serão ajustados no
robô físico por otimização por enxame de partículas (PSO)
\cite{kennedy1995pso}. No PSO, cada partícula representa um conjunto
candidato de ganhos $\mathbf{x}_i = (K_p, K_i, K_d)$. A cada iteração,
sua velocidade e sua posição são atualizadas por
\begin{equation}
	\label{eq:pso}
	\mathbf{v}_i \leftarrow w\,\mathbf{v}_i
		+ c_1 r_1 (\mathbf{p}_i - \mathbf{x}_i)
		+ c_2 r_2 (\mathbf{g} - \mathbf{x}_i), \qquad
	\mathbf{x}_i \leftarrow \mathbf{x}_i + \mathbf{v}_i,
\end{equation}
em que $\mathbf{p}_i$ é a melhor posição já encontrada pela partícula,
$\mathbf{g}$ é a melhor posição encontrada pelo enxame, $w$ é o peso de
inércia, $c_1$ e $c_2$ são coeficientes de aceleração e $r_1$ e $r_2$ são
números aleatórios entre 0 e 1. Cada avaliação consiste em aplicar os
ganhos no robô, comandar uma sequência de velocidades e medir um índice
de desempenho do erro de velocidade, como a integral do erro absoluto.
Cada tentativa deve durar pelo menos cerca de cinco constantes de tempo
do motor, medidas no ensaio de degrau, para que a roda tenha tempo de
chegar ao regime. Para isso, o nó da base receberá os ganhos como
parâmetros do ROS~2, que o otimizador pode trocar sem reiniciar o nó.
